\subsection{Real-World Problem Description}

Signalized intersections are typically retimed
based on when and how severely queues build up
on their approaches. Traffic engineers
commonly study queues using input--output
(cumulative arrival and departure) analysis:
the number of vehicles waiting on an approach
grows when more vehicles arrive than the signal
discharges, and shrinks otherwise
\cite{daganzo1997, hcm2022}.

Modern intersections are instrumented with
vehicle detectors that report, for every short
time bin (for example, one minute), how many
vehicles arrived at the stop line and how many
departed through it. Over days or weeks of
operation this produces a long sequence of
per-minute counts.

A natural question for a traffic engineer is:
\emph{during which continuous period did the
queue on this approach grow the most?}
This period is the worst congestion episode
in the data. Identifying it tells the engineer
when the current signal timing is least able to
serve demand, and is therefore a candidate
period for retiming or coordination
\cite{fhwa2008}.

The period of interest need not align with
fixed boundaries such as the morning peak hour;
it can start and end at any minute. Short
intervals in which the queue partially drains
may lie inside the worst episode, as long as
the overall build-up across the period is
maximal.

This simplified model assumes that detector
counts are accurate, that a single approach is
studied independently, that time is divided
into equal-length bins, and that the queue is
not bounded by physical storage capacity.
The model identifies the worst episode; it does
not by itself prescribe a new signal plan.

\subsection{Mathematical Abstraction}

Let the observation horizon consist of $n$
consecutive minutes. For minute $i$, let
$a_i$ be the number of arrivals and $d_i$ the
number of departures. The net queue change in
minute $i$ is

\[
x_i = a_i - d_i \in \mathbb{Z}.
\]

The problem is therefore abstracted as a
sequence

\[
X = \langle x_1, x_2, \ldots, x_n \rangle,
\]

where $x_i > 0$ means the queue grew and
$x_i < 0$ means the queue drained.

A candidate period is a contiguous index
interval $[l, r]$ with $1 \le l \le r \le n$.
Let

\[
\mathcal{I}_n = \{ [l, r] : 1 \le l \le r \le n \}
\]

be the set of all candidate periods, so
$|\mathcal{I}_n| = n(n+1)/2$.

The total queue build-up over a period is

\[
S(l, r) = \sum_{k=l}^{r} x_k ,
\]

which, under the input--output model, equals
the difference between the queue length at the
end of minute $r$ and at the start of minute $l$.

The goal is to find

\[
[l^*, r^*] \in \arg\max_{[l,r] \in \mathcal{I}_n} S(l, r).
\]

This is the \emph{maximum subarray} problem
\cite{bentley1984, clrs2009}. Periods are
required to be non-empty; if every $x_i$ is
negative, the answer is the single minute with
the smallest drain.
